\documentclass[10pt,a4paper]{amsart}
\usepackage{amsmath,amssymb,amsthm,graphicx}
\usepackage{verbatim}
\newenvironment{reference}{\comment}{\endcomment}
%\newenvironment{reference}{}{}
\newenvironment{slogan}{\comment}{\endcomment}
\newenvironment{history}{\comment}{\endcomment}

% For commutative diagrams we use Xy-pic
\usepackage[all]{xy}

% We use 2cell for 2-commutative diagrams.
\xyoption{2cell}
% Theorem environments.
%
\theoremstyle{plain}
\newtheorem{theorem}[subsection]{Theorem}
\newtheorem{proposition}[subsection]{Proposition}
\newtheorem{lemma}[subsection]{Lemma}

\theoremstyle{definition}
\newtheorem{definition}[subsection]{Definition}
\newtheorem{example}[subsection]{Example}
\newtheorem{exercise}[subsection]{Exercise}
\newtheorem{situation}[subsection]{Situation}

\theoremstyle{remark}
\newtheorem{remark}[subsection]{Remark}
\newtheorem{remarks}[subsection]{Remarks}

\numberwithin{equation}{subsection}

% Macros
%
\def\lim{\mathop{\mathrm{lim}}\nolimits}
\def\colim{\mathop{\mathrm{colim}}\nolimits}
\def\Spec{\mathop{\mathrm{Spec}}}
\def\Hom{\mathop{\mathrm{Hom}}\nolimits}
\def\Ext{\mathop{\mathrm{Ext}}\nolimits}
\def\SheafHom{\mathop{\mathcal{H}\!\mathit{om}}\nolimits}
\def\SheafExt{\mathop{\mathcal{E}\!\mathit{xt}}\nolimits}
\def\Sch{\mathit{Sch}}
\def\Mor{\mathop{\mathrm{Mor}}\nolimits}
\def\Ob{\mathop{\mathrm{Ob}}\nolimits}
\def\Sh{\mathop{\mathit{Sh}}\nolimits}
\def\NL{\mathop{N\!L}\nolimits}
\def\CH{\mathop{\mathrm{CH}}\nolimits}
\def\proetale{{pro\text{-}\acute{e}tale}}
\def\etale{{\acute{e}tale}}
\def\QCoh{\mathit{QCoh}}
\def\Ker{\mathop{\mathrm{Ker}}}
\def\Im{\mathop{\mathrm{Im}}}
\def\Coker{\mathop{\mathrm{Coker}}}
\def\Coim{\mathop{\mathrm{Coim}}}

% Boxtimes
%
\DeclareMathSymbol{\boxtimes}{\mathbin}{AMSa}{"02}

%
% Macros for moduli stacks/spaces
%
\def\QCohstack{\mathcal{QC}\!\mathit{oh}}
\def\Cohstack{\mathcal{C}\!\mathit{oh}}
\def\Spacesstack{\mathcal{S}\!\mathit{paces}}
\def\Quotfunctor{\mathrm{Quot}}
\def\Hilbfunctor{\mathrm{Hilb}}
\def\Curvesstack{\mathcal{C}\!\mathit{urves}}
\def\Polarizedstack{\mathcal{P}\!\mathit{olarized}}
\def\Complexesstack{\mathcal{C}\!\mathit{omplexes}}
% \Pic is the operator that assigns to X its picard group, usage \Pic(X)
% \Picardstack_{X/B} denotes the Picard stack of X over B
% \Picardfunctor_{X/B} denotes the Picard functor of X over B
\def\Pic{\mathop{\mathrm{Pic}}\nolimits}
\def\Picardstack{\mathcal{P}\!\mathit{ic}}
\def\Picardfunctor{\mathrm{Pic}}
\def\Deformationcategory{\mathcal{D}\!\mathit{ef}}
\usepackage{cleveref}
\crefname{originalreference}{original reference}{original references}
\Crefname{originalreference}{Original reference}{Original references}
\makeatletter
\newlabel{situation-dualizing}{{0AU4}{source}{}{sourcepointer-1}{}}
\newlabel{situation-dualizing@cref}{{[originalreference][0][]0AU4}{[1][0][]source}{}{}{}}
\makeatother
\crefname{originalreference}{original reference}{original references}
\Crefname{originalreference}{Original reference}{Original references}
\makeatletter
\newlabel{item-cocycle-glueing}{{0AU6}{source}{}{sourcepointer-2}{}}
\newlabel{item-cocycle-glueing@cref}{{[originalreference][0][]0AU6}{[1][0][]source}{}{}{}}
\makeatother
\crefname{originallemma}{original lemma}{original lemmas}
\Crefname{originallemma}{Original lemma}{Original lemmas}
\makeatletter
\newlabel{lemma-good-dualizing-unique}{{0AU7}{source}{}{sourcepointer-3}{}}
\newlabel{lemma-good-dualizing-unique@cref}{{[originallemma][0][]0AU7}{[1][0][]source}{}{}{}}
\makeatother
\title{The normalised dualizing complex is independent of the chosen finite separated open covering}
\author{The Stacks Project Authors}
\date{stacks-project:duality.tex:local-snapshot-20261012; GNU FDL 1.2 or later}
\begin{document}
\maketitle
\noindent\textit{Editorial note.} All delivered mathematical text is the author's own, copied from the pinned Stacks Project chapter file duality.tex; no statement or proof step was rewritten, reordered or removed.\par\medskip
\noindent\textit{Editorial note.} Preamble provenance: the pinned chapter file opens with an input of the shared preamble and carries no preamble text of its own, so the project's single Stacks-wide preamble staged at sources/stacks-duality/preamble.tex is excerpted verbatim under the recipe fixed by LANE\_BRIEF 22.1, as three spans, lines 8-13 (verbatim package and the reference/slogan/history environments), lines 14-18 (Xy-pic setup) and lines 154-217 (the native theorem declarations and mathematical macros). Nothing was added, removed or redefined; the wrapper class and the amsmath/amssymb/amsthm/graphicx base packages are supplied by the assembler.\par\medskip
\noindent\textit{Editorial note.} The counted statement and proof contain no cross reference, no citation and no included file, so the delivered body is self contained.\par\medskip
\noindent\textit{Editorial note.} Canonical location of the counted proposition: Stacks tag 0AU8 (sources/stacks-duality/pinned-original-source.tex line 5348).\par\medskip
\noindent\textit{Allowed standard prerequisites (editorial).}\begin{itemize}
\item The situation, the cocycle gluing condition and the notion of a normalised dualizing complex. Same chapter, Situation labelled situation-dualizing and equation labelled item-cocycle-glueing, invoked in the statement and cited at proof lines 5378-5384. Fixes the data and the compatibility condition.
\item Uniqueness of a normalised dualizing complex relative to a covering. Same chapter, Lemma labelled lemma-good-dualizing-unique, cited at proof lines 5185-5189 as referenced in the proof text. Produces the unique compatibility isomorphism alpha\_\{n+1\} on the new open.
\end{itemize}\medskip
\ifdefined\hypertarget\hypertarget{sourcepointer-1}{}\fi
\noindent\textit{Source reference.} Original reference 0AU4: sources/stacks-duality/pinned-original-source.tex, source lines 5259--5259. Same-chapter pointer into the author environment situation at source line 5259; invoked at proof line None of the counted proof. It is retained only as a navigation pointer; every load-bearing step is inside the delivered span.\par\medskip
\ifdefined\hypertarget\hypertarget{sourcepointer-2}{}\fi
\noindent\textit{Source reference.} Original reference 0AU6: sources/stacks-duality/pinned-original-source.tex, source lines 5285--5285. Same-chapter pointer into the author environment enumerate at source line 5285; invoked at proof line 5377 of the counted proof. It is retained only as a navigation pointer; every load-bearing step is inside the delivered span.\par\medskip
\ifdefined\hypertarget\hypertarget{sourcepointer-3}{}\fi
\noindent\textit{Source reference.} Original lemma 0AU7: sources/stacks-duality/pinned-original-source.tex, source lines 5317--5317. Same-chapter pointer into the author environment lemma at source line 5317; invoked at proof line 5382 of the counted proof. It is retained only as a navigation pointer; every load-bearing step is inside the delivered span.\par\medskip
\begin{lemma}
\label{lemma-good-dualizing-independence-covering}
In Situation \ref{situation-dualizing} let $X$ be a scheme of finite type
over $S$ and let $\mathcal{U}$, $\mathcal{V}$ be two finite open coverings
of $X$ by schemes separated over $S$.
If there exists a dualizing complex normalized
relative to $\omega_S^\bullet$ and $\mathcal{U}$, then
there exists a dualizing complex normalized relative to
$\omega_S^\bullet$ and $\mathcal{V}$ and these complexes are
canonically isomorphic.
\end{lemma}

\begin{proof}
It suffices to prove this when $\mathcal{U}$ is given by the opens
$U_1, \ldots, U_n$ and $\mathcal{V}$ by the opens $U_1, \ldots, U_{n + m}$.
In fact, we may and do even assume $m = 1$.
To go from a dualizing complex $(K, \alpha_i)$ normalized
relative to $\omega_S^\bullet$ and $\mathcal{V}$ to a
dualizing complex normalized relative to $\omega_S^\bullet$ and $\mathcal{U}$
is achieved by forgetting about $\alpha_i$ for $i = n + 1$. Conversely, let
$(K, \alpha_i)$ be a dualizing complex normalized relative to
$\omega_S^\bullet$ and $\mathcal{U}$.
To finish the proof we need to construct a map
$\alpha_{n + 1} : K|_{U_{n + 1}} \to \omega_{n + 1}^\bullet$ satisfying
the desired conditions.
To do this we observe that $U_{n + 1} = \bigcup U_i \cap U_{n + 1}$
is an open covering.
It is clear that $(K|_{U_{n + 1}}, \alpha_i|_{U_i \cap U_{n + 1}})$
is a dualizing complex normalized relative to $\omega_S^\bullet$
and the covering $U_{n + 1} = \bigcup U_i \cap U_{n + 1}$.
On the other hand, by condition (\ref{item-cocycle-glueing}) the pair
$(\omega_{n + 1}^\bullet|_{U_{n + 1}}, \varphi_{n + 1i})$
is another dualizing complex normalized relative to $\omega_S^\bullet$
and the covering
$U_{n + 1} = \bigcup U_i \cap U_{n + 1}$.
By Lemma \ref{lemma-good-dualizing-unique} we obtain a unique isomorphism
$$
\alpha_{n + 1} : K|_{U_{n + 1}} \longrightarrow \omega_{n + 1}^\bullet
$$
compatible with the given local isomorphisms.
It is a pleasant exercise to show that this means it satisfies
the required property.
\end{proof}

\end{document}
